\chapter{Nano- and Microfabrication}\label{sec:fabrication}

This chapter outlines the fabrication processes central to this work, beginning with the motivation behind structuring diamond at the nanoscale. We provide a comparative analysis of current scanning probe architectures, highlighting the evolution from simple vertical pillars to parabolic geometries. We then introduce the structures developed in this thesis as an evolutionary step forward, designed to enhance sensitivity through the use an ensemble of NVs and while retaining the high spatial resolution required for scanning applications. Detailed discussions on design choices, mask materials, and specific microfabrication steps are presented to offer a complete view of the considerations in developing a scanning probe.

\section{Introduction}
Color centers in diamond is a leading and continuously developing platform for many quantum technologies in quantum information applications\cite{childressDiamondNVCenters2013, wrachtrupProcessingQuantumInformation2006} and advanced sensing\cite{taylorHighsensitivityDiamondMagnetometer2008, schirhaglNitrogenVacancyCentersDiamond2014}, by virtue of its excellent spin properties and optical addressability. The utility of the NV center and other color centers is fundamentally dependant on the properties of the host material, diamond. One one hand, diamond has many advantages such as a wide electronic bandgap of $\sim$\qty{5.5}{\electronvolt}, high thermal conductivity, mechanical robustness, chemical inertness, and bio-compatibility making it an attractive choice for many applications. Crystalline defects in diamond in some combination of displacement of substitution of native carbon exhibits many optically active defects, electronically isolated in the bulk band gap, and with control over spatial density, entering diamond in the realm of quantum optics. These defects also exhibit low phonon-induced relaxation due to the high Debye temperature of diamond. Furthermore, the availability of single crystal diamonds, with a relatively low natural abundance of nuclear spins leads to a magnetic noise free environment making it a perfect host for spins with long coherence times\cite{balasubramanianUltralongSpinCoherence2009}.

The conventional and convenient optical addressability of defects in diamond means efficiently collecting the emitted photons; due to a high refractive index of $\sim$2.4 and a resulting total internal reflection critical angle of $\sim$\qty{25}{\degree}\cite{zhuMulticoneDiamondWaveguides2023a}; is challenging. In this regard, diamond nanofabrication is an enabling development for many quantum technologies and thus emerges the field of diamond nanophotonics\cite{andoSmoothHighrateReactive2002, baynTriangularNanobeamPhotonic2011}.

\section{Nanophotonics and quantum sensing}

The field of quantum sensing has made great advances making use of the ultra long coherence times approaching a few \qty{}{ms}\cite{balasubramanianUltralongSpinCoherence2009} and magnetic mapping with nanometer spatial resolution\cite{appelFabricationAllDiamond2016} possible with NV centers in diamond. This has established magnetic field sensing with high sensitivity and high spatial resolution as one of the key applications of NV centers, even seeing commercial viability. An ideal magnetometer would be able to approach infinitesimally close to the magnetic field source and probe the magnetic signal for infinitely long, and convey this information via a measurement with an infinite signal-to-noise(SNR) ratio. To approach these ideals for a NV based magnetometer, it is desirable to then to employ near surface NV centers, to achieve close proximity to the source; to retain good coherence properties of the NV center, to enable long probing times; and to fabricate a nanophotonic waveguide that enhances the collected fluorescence, increasing the sensitivity.